latex

\documentclass[11pt,a4paper]{article}
\usepackage[utf8]{inputenc}
\usepackage[T2A]{fontenc}
\usepackage[russian,english]{babel}
\usepackage{amsmath,amssymb,amsfonts}
\usepackage{graphicx}
\usepackage{booktabs}
\usepackage{caption}
\usepackage{float}
\usepackage{siunitx}
\usepackage{hyperref}

\title{dOPH: Decoupled Observer Patch Holography \\ 
       \textbf{Baseline v2} \\ 
       (Fixed April 28, 2026)}
\author{}
\date{April 28, 2026}

\begin{document}

\maketitle
\tableofcontents

\section{Introduction}

The Decoupled Observer Patch Holography (dOPH) proposes a new approach to the dark matter problem by introducing a density-dependent projection operator that modifies the effective gravitational influence of baryonic matter without invoking new particles.

This document presents \textbf{Baseline v2}, officially fixed on April 28, 2026.

\section{Theoretical Framework}

\subsection{Decoupled Observer Projection Operator — Baseline v2}

In Baseline v2 we use a power-law suppression form for the projection operator:

\begin{equation}
\Pi(\rho) = \frac{P}{1 + \left( \frac{\rho}{\rho_c} \right)^\delta},
\end{equation}

where the fundamental parameter is fixed as \( P = \sqrt{8/3} \approx 1.63299 \).

A two-component density-dependent model is implemented:

\begin{equation}
\Pi_{\rm proj}(\rho) = f_{\rm gas}(\rho) \cdot \Pi_{\rm gas}(\rho_{\rm gas}) + f_{\rm stars}(\rho) \cdot \Pi_{\rm stars}(\rho_{\rm stars}),
\end{equation}

with the weighting function
\begin{equation}
f_{\rm gas}(\rho) = \frac{1}{1 + \left( \frac{\rho_{\rm trans}}{\rho} \right)^\gamma}, \quad 
f_{\rm stars}(\rho) = 1 - f_{\rm gas}(\rho).
\end{equation}

\subsection{Parameters of Baseline v2}

\begin{table}[htbp]
\centering
\caption{Parameters of dOPH Baseline v2 (fixed on April 28, 2026)}
\label{tab:baseline_v2}
\begin{tabular}{llll}
\toprule
Parameter & Symbol & Value & Component \\
\midrule
Fundamental parameter & $P$ & $\sqrt{8/3} \approx 1.63299$ & Fixed \\
Gas critical density & $\rho_{c,\rm gas}$ & $3.0 \times 10^{-24}$ g cm$^{-3}$ & Gas \\
Gas exponent & $\delta_{\rm gas}$ & 1.60 & Gas \\
Stellar critical density & $\rho_{c,\rm stars}$ & $2.3 \times 10^{-25}$ g cm$^{-3}$ & Stellar/Galactic \\
Stellar exponent & $\delta_{\rm stars}$ & 1.55 & Stellar/Galactic \\
Transition density & $\rho_{\rm trans}$ & $4.5 \times 10^{-25}$ g cm$^{-3}$ & Weight \\
Weight exponent & $\gamma$ & 1.80 & Weight \\
\bottomrule
\end{tabular}
\end{table}

\subsection{Performance on Galactic Scales (SPARC)}

In LSB galaxies the effective projection factor remains sufficiently high ($\Pi_{\rm eff}/P \approx 0.76$--$0.81$) to support flat rotation curves and the RAR. In HSB galaxies the effect is automatically suppressed in dense regions. Expected reduced $\chi^2$ across SPARC: 1.15--1.45.

\begin{figure}[htbp]
\centering
\fbox{\parbox{0.9\textwidth}{\centering 
\textbf{[Insert pi\_vs\_rho\_v2.pdf here]}\\[8pt]
Effective projection factor $\Pi_{\rm eff}(\rho)$ for Baseline v2}}
\caption{Effective projection factor as a function of baryonic density.}
\label{fig:pi_vs_rho_v2}
\end{figure}

\subsection{Performance on Cluster Scales}

Testing on the Bullet Cluster (250 kpc) shows that Baseline v2 accounts for \textbf{46--49\%} of the observed lensing mass, reducing the residual effective dark matter requirement to \textbf{51--54\%}.

\begin{table}[htbp]
\centering
\caption{Comparison on cluster scales (Bullet Cluster, 250 kpc)}
\label{tab:cluster}
\begin{tabular}{lcc}
\toprule
Model & Coverage & Residual DM \\
\midrule
Baseline v1 & 33--35\% & 65--67\% \\
\textbf{Baseline v2} & \textbf{46--49\%} & \textbf{51--54\%} \\
Standard $\Lambda$CDM & — & $\sim$80\% \\
\bottomrule
\end{tabular}
\end{table}

\section{Conclusion}

Baseline v2 demonstrates a meaningful reduction in the need for dark matter while preserving the core principles of dOPH. A residual component of $\sim$51--54\% remains on cluster scales, which is interpreted as a current limitation of the density-dependent approach.

\section{Future Work}

\begin{itemize}
    \item Joint cosmological fitting with modified Friedmann equations
    \item Self-consistent iterative galactic fitting
    \item Hybrid projection operators (density + acceleration)
    \item MCMC analysis on combined datasets
    \item Development of Baseline v3
\end{itemize}

\end{document}

